%% ODER: format ==         = "\mathrel{==}"
%% ODER: format /=         = "\neq "
%
%
\makeatletter
\@ifundefined{lhs2tex.lhs2tex.sty.read}%
  {\@namedef{lhs2tex.lhs2tex.sty.read}{}%
   \newcommand\SkipToFmtEnd{}%
   \newcommand\EndFmtInput{}%
   \long\def\SkipToFmtEnd#1\EndFmtInput{}%
  }\SkipToFmtEnd

\newcommand\ReadOnlyOnce[1]{\@ifundefined{#1}{\@namedef{#1}{}}\SkipToFmtEnd}
\usepackage{amstext}
\usepackage{amssymb}
\usepackage{stmaryrd}
\DeclareFontFamily{OT1}{cmtex}{}
\DeclareFontShape{OT1}{cmtex}{m}{n}
  {<5><6><7><8>cmtex8
   <9>cmtex9
   <10><10.95><12><14.4><17.28><20.74><24.88>cmtex10}{}
\DeclareFontShape{OT1}{cmtex}{m}{it}
  {<-> ssub * cmtt/m/it}{}
\newcommand{\texfamily}{\fontfamily{cmtex}\selectfont}
\DeclareFontShape{OT1}{cmtt}{bx}{n}
  {<5><6><7><8>cmtt8
   <9>cmbtt9
   <10><10.95><12><14.4><17.28><20.74><24.88>cmbtt10}{}
\DeclareFontShape{OT1}{cmtex}{bx}{n}
  {<-> ssub * cmtt/bx/n}{}
\newcommand{\tex}[1]{\text{\texfamily#1}}	% NEU

\newcommand{\Sp}{\hskip.33334em\relax}


\newcommand{\Conid}[1]{\mathit{#1}}
\newcommand{\Varid}[1]{\mathit{#1}}
\newcommand{\anonymous}{\kern0.06em \vbox{\hrule\@width.5em}}
\newcommand{\plus}{\mathbin{+\!\!\!+}}
\newcommand{\bind}{\mathbin{>\!\!\!>\mkern-6.7mu=}}
\newcommand{\rbind}{\mathbin{=\mkern-6.7mu<\!\!\!<}}% suggested by Neil Mitchell
\newcommand{\sequ}{\mathbin{>\!\!\!>}}
\renewcommand{\leq}{\leqslant}
\renewcommand{\geq}{\geqslant}
\usepackage{polytable}

%mathindent has to be defined
\@ifundefined{mathindent}%
  {\newdimen\mathindent\mathindent\leftmargini}%
  {}%

\def\resethooks{%
  \global\let\SaveRestoreHook\empty
  \global\let\ColumnHook\empty}
\newcommand*{\savecolumns}[1][default]%
  {\g@addto@macro\SaveRestoreHook{\savecolumns[#1]}}
\newcommand*{\restorecolumns}[1][default]%
  {\g@addto@macro\SaveRestoreHook{\restorecolumns[#1]}}
\newcommand*{\aligncolumn}[2]%
  {\g@addto@macro\ColumnHook{\column{#1}{#2}}}

\resethooks

\newcommand{\onelinecommentchars}{\quad-{}- }
\newcommand{\commentbeginchars}{\enskip\{-}
\newcommand{\commentendchars}{-\}\enskip}

\newcommand{\visiblecomments}{%
  \let\onelinecomment=\onelinecommentchars
  \let\commentbegin=\commentbeginchars
  \let\commentend=\commentendchars}

\newcommand{\invisiblecomments}{%
  \let\onelinecomment=\empty
  \let\commentbegin=\empty
  \let\commentend=\empty}

\visiblecomments

\newlength{\blanklineskip}
\setlength{\blanklineskip}{0.66084ex}

\newcommand{\hsindent}[1]{\quad}% default is fixed indentation
\let\hspre\empty
\let\hspost\empty
\newcommand{\NB}{\textbf{NB}}
\newcommand{\Todo}[1]{$\langle$\textbf{To do:}~#1$\rangle$}

\EndFmtInput
\makeatother
%
%
%
%
%
%
% This package provides two environments suitable to take the place
% of hscode, called "plainhscode" and "arrayhscode". 
%
% The plain environment surrounds each code block by vertical space,
% and it uses \abovedisplayskip and \belowdisplayskip to get spacing
% similar to formulas. Note that if these dimensions are changed,
% the spacing around displayed math formulas changes as well.
% All code is indented using \leftskip.
%
% Changed 19.08.2004 to reflect changes in colorcode. Should work with
% CodeGroup.sty.
%
\ReadOnlyOnce{polycode.fmt}%
\makeatletter

\newcommand{\hsnewpar}[1]%
  {{\parskip=0pt\parindent=0pt\par\vskip #1\noindent}}

% can be used, for instance, to redefine the code size, by setting the
% command to \small or something alike
\newcommand{\hscodestyle}{}

% The command \sethscode can be used to switch the code formatting
% behaviour by mapping the hscode environment in the subst directive
% to a new LaTeX environment.

\newcommand{\sethscode}[1]%
  {\expandafter\let\expandafter\hscode\csname #1\endcsname
   \expandafter\let\expandafter\endhscode\csname end#1\endcsname}

% "compatibility" mode restores the non-polycode.fmt layout.

\newenvironment{compathscode}%
  {\par\noindent
   \advance\leftskip\mathindent
   \hscodestyle
   \let\\=\@normalcr
   \let\hspre\(\let\hspost\)%
   \pboxed}%
  {\endpboxed\)%
   \par\noindent
   \ignorespacesafterend}

\newcommand{\compaths}{\sethscode{compathscode}}

% "plain" mode is the proposed default.
% It should now work with \centering.
% This required some changes. The old version
% is still available for reference as oldplainhscode.

\newenvironment{plainhscode}%
  {\hsnewpar\abovedisplayskip
   \advance\leftskip\mathindent
   \hscodestyle
   \let\hspre\(\let\hspost\)%
   \pboxed}%
  {\endpboxed%
   \hsnewpar\belowdisplayskip
   \ignorespacesafterend}

\newenvironment{oldplainhscode}%
  {\hsnewpar\abovedisplayskip
   \advance\leftskip\mathindent
   \hscodestyle
   \let\\=\@normalcr
   \(\pboxed}%
  {\endpboxed\)%
   \hsnewpar\belowdisplayskip
   \ignorespacesafterend}

% Here, we make plainhscode the default environment.

\newcommand{\plainhs}{\sethscode{plainhscode}}
\newcommand{\oldplainhs}{\sethscode{oldplainhscode}}
\plainhs

% The arrayhscode is like plain, but makes use of polytable's
% parray environment which disallows page breaks in code blocks.

\newenvironment{arrayhscode}%
  {\hsnewpar\abovedisplayskip
   \advance\leftskip\mathindent
   \hscodestyle
   \let\\=\@normalcr
   \(\parray}%
  {\endparray\)%
   \hsnewpar\belowdisplayskip
   \ignorespacesafterend}

\newcommand{\arrayhs}{\sethscode{arrayhscode}}

% The mathhscode environment also makes use of polytable's parray 
% environment. It is supposed to be used only inside math mode 
% (I used it to typeset the type rules in my thesis).

\newenvironment{mathhscode}%
  {\parray}{\endparray}

\newcommand{\mathhs}{\sethscode{mathhscode}}

% texths is similar to mathhs, but works in text mode.

\newenvironment{texthscode}%
  {\(\parray}{\endparray\)}

\newcommand{\texths}{\sethscode{texthscode}}

% The framed environment places code in a framed box.

\def\codeframewidth{\arrayrulewidth}
\RequirePackage{calc}

\newenvironment{framedhscode}%
  {\parskip=\abovedisplayskip\par\noindent
   \hscodestyle
   \arrayrulewidth=\codeframewidth
   \tabular{@{}|p{\linewidth-2\arraycolsep-2\arrayrulewidth-2pt}|@{}}%
   \hline\framedhslinecorrect\\{-1.5ex}%
   \let\endoflinesave=\\
   \let\\=\@normalcr
   \(\pboxed}%
  {\endpboxed\)%
   \framedhslinecorrect\endoflinesave{.5ex}\hline
   \endtabular
   \parskip=\belowdisplayskip\par\noindent
   \ignorespacesafterend}

\newcommand{\framedhslinecorrect}[2]%
  {#1[#2]}

\newcommand{\framedhs}{\sethscode{framedhscode}}

% The inlinehscode environment is an experimental environment
% that can be used to typeset displayed code inline.

\newenvironment{inlinehscode}%
  {\(\def\column##1##2{}%
   \let\>\undefined\let\<\undefined\let\\\undefined
   \newcommand\>[1][]{}\newcommand\<[1][]{}\newcommand\\[1][]{}%
   \def\fromto##1##2##3{##3}%
   \def\nextline{}}{\) }%

\newcommand{\inlinehs}{\sethscode{inlinehscode}}

% The joincode environment is a separate environment that
% can be used to surround and thereby connect multiple code
% blocks.

\newenvironment{joincode}%
  {\let\orighscode=\hscode
   \let\origendhscode=\endhscode
   \def\endhscode{\def\hscode{\endgroup\def\@currenvir{hscode}\\}\begingroup}
   %\let\SaveRestoreHook=\empty
   %\let\ColumnHook=\empty
   %\let\resethooks=\empty
   \orighscode\def\hscode{\endgroup\def\@currenvir{hscode}}}%
  {\origendhscode
   \global\let\hscode=\orighscode
   \global\let\endhscode=\origendhscode}%

\makeatother
\EndFmtInput
%
%% ----------------------
%% Stuff for TimCore


%% 'g' inferences




%% --------- Effects ------------


%% ----------------------




%% ----------- Applications --------------
%% %format applyx a b = a "@" b
%% %format applyx a b = a "\," b

%% Function application with no arguments

%% Function application with 1 argument

%% Function application with 1 argument, without parenthesis

%% Function application with many arguments

%% ----------- End of Applications --------------


%% Sequences e1; ...; en; v

%% xs etc stole syntax used in examples.

%% Old version: | for BNF, [] for choice
%%%format choice = "[ \! ]"
%%%format `choice` = "\mathop{[ \! ]}"
%% 1mm is about 0.3em with the current font

%% New version: tall | for BNF, fat | for choice
%%format vbar = "\mathop{\bigm|}"
%%format choice = "\boldsymbol{|}"
%%format `choice` = "\mathop{\boldsymbol{|}}"




%% Arrays and tuples
%% %format arrKW = "\textbf{arr}"
%% %format arr (x) = arrKW "\{" x "\}"





%%  OLD SYNTAX   %format split (x) (y) (z) = "\textbf{split}(" x ")\{" y "," z "\}"
%% %format defKW = "\textbf{def}"

% only used in comments, but still prettier this way:
%%%%format map = "\textbf{map}"

%% Effects checking, like succeeds{e}


%% Unification
%%%format == = "\!=\!"


% format := = "\mapsto"
%% %format squiggt = "\!\leadsto\!"
%  format effs x = "\!\!<\!\!" x "\!\!>\!\!" dots


































%% ----------- Metatheory --------------
%%Formatting for skew confluence proof


%%%%%%%%%%%%%%%%%%%%%%%%%%%%%%%%%%%%

\section{Three Languages}

\subsection{Lambda Calculus}

TBD

\subsection{Lambda-Letrec Calculus}

TBD

\subsection{Verse Calculus}


\newcommand{\figureSkewProofRewriteRulesOLD}{
  \begin{figure}
    $\begin{array}{@{}l}
\begin{array}{@{\hskip0.5em} l @{\hspace{-.5em}} r @{\;}c@{\;} l l@{}}
\rulegroupname{Application:}
\rulename{app-add}          & \ensuremath{\textbf{add}\langle\Varid{k}_{\mathrm{1}},\Varid{k}_{\mathrm{2}}\rangle} & \movesto & \ensuremath{\Varid{k}_{\mathrm{3}}} & \text{where $k_3 = k_1 + k_2$}\\
\rulename{app-gt}           & \ensuremath{\textbf{gt}\langle\Varid{k}_{\mathrm{1}},\Varid{k}_{\mathrm{2}}\rangle}   & \movesto & \ensuremath{\Varid{k}_{\mathrm{1}}} &  \text{if $k_1 > k_2$} \\
\rulename{app-gt-fail}      & \ensuremath{\textbf{gt}\langle\Varid{k}_{\mathrm{1}},\Varid{k}_{\mathrm{2}}\rangle}   & \movesto & \ensuremath{\textbf{fail}} &  \text{if $k_1 \leq k_2$} \\
\rulename{app-beta}\rlap{${}^{\alpha}$} & \ensuremath{(\lambda\Varid{x}.\,\Varid{e})(\Varid{v})} & \movesto & \ensuremath{\exists\Varid{x}.\,\Varid{x}\hspace{-0.14em}=\hspace{-0.14em}\Varid{v};\,\Varid{e}}
                                                  & \text{if $x \not\in \freevars{\ensuremath{\Varid{v}}}$} \\
\rulename{app-tup}          & \ensuremath{\langle{}\Varid{v}_{\mathrm{0}},\mydots,\Varid{v}_{\Varid{n}}\rangle(\Varid{v})} & \movesto &
                              \multicolumn{2}{@{}l}{\ensuremath{\exists\Varid{x}.\,\Varid{x}\hspace{-0.14em}=\hspace{-0.14em}\Varid{v};(\Varid{x}\hspace{-0.14em}=\hspace{-0.14em}\mathrm{0};\Varid{v}_{\mathrm{0}})\;\rule[-0.1em]{0.15em}{0.75em}\;\mydots\;\rule[-0.1em]{0.15em}{0.75em}\;(\Varid{x}\hspace{-0.14em}=\hspace{-0.14em}\Varid{n};\Varid{v}_{\Varid{n}})}
                              \hspace{0.7em}\text{fresh $\ensuremath{\Varid{x}} \not\in \freevars{\ensuremath{\Varid{v},\Varid{v}_{\mathrm{0}},\mydots,\Varid{v}_{\Varid{n}}}}$}} \hskip-0.7em \\
\rulename{app-tup-0}        & \ensuremath{\langle\rangle(\Varid{v})} & \movesto & \ensuremath{\textbf{fail}} \\
\\[-0.6em]
\rulegroupname{Unification:}
  \rulename{u-lit}   & \ensuremath{\Varid{k}_{\mathrm{1}}\hspace{-0.14em}=\hspace{-0.14em}\Varid{k}_{\mathrm{2}};\,\Varid{e}} & \movesto & \ensuremath{\Varid{e}} & \text{if $k_1=k_2$} \\
  \rulename{u-tup} & \hspace{-2.0em} \ensuremath{\langle{}\Varid{v}_{\mathrm{1}},\mydots,\Varid{v}_{\Varid{n}}\rangle\hspace{-0.14em}=\hspace{-0.14em}\langle{}\Varid{v}_{1}',\mydots,\Varid{v}_{\Varid{n}}'\rangle;\,\Varid{e}} & \movesto & \ensuremath{\Varid{v}_{\mathrm{1}}\hspace{-0.14em}=\hspace{-0.14em}\Varid{v}_{1}';\,\mydots;\,\Varid{v}_{\Varid{n}}\hspace{-0.14em}=\hspace{-0.14em}\Varid{v}_{\Varid{n}}';\,\Varid{e}} \\
  \rulename{u-fail} & \ensuremath{\mathit{hnf}_{\!1}\hspace{-0.14em}=\hspace{-0.14em}\mathit{hnf}_{\!2};\,\Varid{e}} & \movesto & \multicolumn{2}{l}{\ensuremath{\textbf{fail}}
       \hspace{2cm} \begin{array}[t]{l} \text{if \rulename{u-lit}, \rulename{u-tup} do not match} \\
         \text{and neither \ensuremath{\mathit{hnf}_{\!1}} nor \ensuremath{\mathit{hnf}_{\!2}} is a lambda}
         \end{array}} \\
  \rulename{u-occurs} & \ensuremath{\Varid{x}\hspace{-0.14em}=\hspace{-0.14em}\Conid{V} [\mskip1.5mu \Varid{x}\mskip1.5mu];\,\Varid{e}} & \movesto & \ensuremath{\textbf{fail}} & \text{if $\ensuremath{\Conid{V}} \ne \ensuremath{\hole}$} \\
  \rulename{u-wrap} & \ensuremath{\Varid{x}\mathrel{=}\mathit{hnf};\,\Varid{e}} & \movesto & \ensuremath{\Varid{x}\mathrel{=}\mu\Varid{x}.\,\mathit{hnf};\,\Varid{e}} &
                      \multicolumn{1}{@{}l@{}}{\begin{array}[t]{l} \text{if \rulename{u-occurs} does not match and} \\
         \text{~~$\ensuremath{\Varid{x}} \in \freevars{\ensuremath{\mathit{hnf}}}$}
         \end{array}} \\
  \rulename{u-unwrap} & \ensuremath{\mu\Varid{x}.\,\mathit{hnf}} & \movesto & \ensuremath{\subst{\mathit{hnf}}{\Varid{x}}{\mu\Varid{x}.\,\mathit{hnf}}} \\
  \rulename{subst} & \ensuremath{\Varid{x}\hspace{-0.14em}=\hspace{-0.14em}\Varid{v};\,\Varid{e}} & \movesto & \ensuremath{\Varid{x}\hspace{-0.14em}=\hspace{-0.14em}\Varid{v};\,\subst{\Varid{e}}{\Varid{x}}{\Varid{v}}} \hspace*{-0.3em} & \text{if $\ensuremath{\Varid{x}} \notin \freevars{\ensuremath{\Varid{v}}}$} \\
  \rulename{hnf-swap} & \ensuremath{\mathit{hnf}\hspace{-0.14em}=\hspace{-0.14em}\Varid{x};\,\Varid{e}}    & \movesto & \ensuremath{\Varid{x}\hspace{-0.14em}=\hspace{-0.14em}\mathit{hnf};\,\Varid{e}} \\
  \rulename{var-swap} & \ensuremath{\Varid{y}\hspace{-0.14em}=\hspace{-0.14em}\Varid{x};\,\Varid{e}}    & \movesto & \ensuremath{\Varid{x}\hspace{-0.14em}=\hspace{-0.14em}\Varid{y};\,\Varid{e}} \\
  \rulename{seq-swap} & \ensuremath{eq;\,\Varid{x}\hspace{-0.14em}=\hspace{-0.14em}\Varid{v};\,\Varid{e}} & \movesto & \ensuremath{\Varid{x}\hspace{-0.14em}=\hspace{-0.14em}\Varid{v};\,eq;\,\Varid{e}} \\

    \\[-0.6em]

\rulegroupname{Elimination:}
\rulename{val-elim} & \ensuremath{\Varid{v};\,\Varid{e}} & \movesto & \ensuremath{\Varid{e}} \\
\rulename{exi-elim} & \ensuremath{\exists\Varid{x}.\,\Varid{e}} & \movesto & \ensuremath{\Varid{e}} & \text{if $\ensuremath{\Varid{x}} \notin \freevars{\ensuremath{\Varid{e}}}$}\\
  \rulename{eqn-elim} & \ensuremath{\exists\Varid{x}.\,\Varid{x}\hspace{-0.14em}=\hspace{-0.14em}\Varid{v};\,\Varid{e}} & \movesto & \ensuremath{\Varid{e}} & \text{if $\ensuremath{\Varid{x}} \notin \freevars{\ensuremath{\Varid{v}},\ensuremath{\Varid{e}}}$} \\
\rulename{fail-elim-eq} & \ensuremath{\Varid{x}\mathrel{=}\textbf{fail};\,\Varid{e}} & \movesto & \ensuremath{\textbf{fail}} & \\
\rulename{fail-elim-l} & \ensuremath{\textbf{fail};\,\Varid{e}} & \movesto & \ensuremath{\textbf{fail}} & \\
\rulename{fail-elim-r} & \ensuremath{eq;\,\textbf{fail}} & \movesto & \ensuremath{\textbf{fail}} & \\

    \\[-0.6em]

\rulegroupname{Normalization:}
\rulename{exi-float-eq}\rlap{${}^{\alpha}$} & \ensuremath{\Varid{v}\mathrel{=}(\exists\Varid{x}.\,\Varid{e});\,\Varid{e'}} & \movesto & \ensuremath{\exists\Varid{x}.\,\Varid{v}\mathrel{=}\Varid{e};\,\Varid{e'}} & \text{if $x \notin \freevars{\ensuremath{\Varid{v}},\ensuremath{\Varid{e'}}}$} \\
\rulename{exi-float-l}\rlap{${}^{\alpha}$} & \ensuremath{(\exists\Varid{x}.\,\Varid{e});\,\Varid{e'}} & \movesto & \ensuremath{\exists\Varid{x}.\,\Varid{e};\,\Varid{e'}} & \text{if $x \notin \freevars{\ensuremath{\Varid{e'}}}$} \\
\rulename{exi-float-r}\rlap{${}^{\alpha}$} & \ensuremath{eq;\,(\exists\Varid{x}.\,\Varid{e})} & \movesto & \ensuremath{\exists\Varid{x}.\,eq;\,\Varid{e}} & \text{if $x \notin \freevars{\ensuremath{eq}}$} \\
\rulename{seq-assoc} & \ensuremath{(eq;\,\Varid{e}_{\mathrm{1}});\,\Varid{e}_{\mathrm{2}}} & \movesto & \ensuremath{eq;\,(\Varid{e}_{\mathrm{1}};\,\Varid{e}_{\mathrm{2}})} \\
\rulename{eqn-float} & \ensuremath{\Varid{v}\hspace{-0.14em}=\hspace{-0.14em}(eq;\,\Varid{e}_{\mathrm{1}});\,\Varid{e}_{\mathrm{2}}} & \movesto & \ensuremath{eq;\,(\Varid{v}\hspace{-0.14em}=\hspace{-0.14em}\Varid{e}_{\mathrm{1}};\,\Varid{e}_{\mathrm{2}})} \\
\rulename{exi-swap} & \ensuremath{\exists\Varid{x}.\,\exists\Varid{y}.\,\Varid{e}} & \movesto & \ensuremath{\exists\Varid{y}.\,\exists\Varid{x}.\,\Varid{e}} \\

    \\[-0.6em]

\rulegroupname{Choice:}
\rulename{one-fail}   & \ensuremath{\textbf{one}\{\textbf{fail}\}} & \movesto & \ensuremath{\textbf{fail}} \\
\rulename{one-value}  & \ensuremath{\textbf{one}\{\Varid{v}\}} & \movesto & \ensuremath{\Varid{v}} \\
\rulename{one-choice} & \ensuremath{\textbf{one}\{\Varid{v}\hspace{0.2em}\mathop{\rule[-0.1em]{0.15em}{0.75em}}\hspace{0.2em}\Varid{e}\}} & \movesto & \ensuremath{\Varid{v}} \\
\rulename{all-fail}   & \ensuremath{\textbf{all}\{\textbf{fail}\}} & \movesto & \ensuremath{\langle{}\rangle} \\
\rulename{all-value}  & \ensuremath{\textbf{all}\{\Varid{v}\}} & \movesto & \ensuremath{\langle{}\Varid{v}\rangle} \\
\rulename{all-choice} & \ensuremath{\textbf{all}\{\Varid{v}_{\mathrm{1}}\;\rule[-0.1em]{0.15em}{0.75em}\;\mydots\;\rule[-0.1em]{0.15em}{0.75em}\;\Varid{v}_{\Varid{n}}\}} & \movesto & \ensuremath{\langle{}\Varid{v}_{\mathrm{1}},\mydots,\Varid{v}_{\Varid{n}}\rangle} \\
\rulename{choose-r} & \ensuremath{\textbf{fail}\;\rule[-0.1em]{0.15em}{0.75em}\;\Varid{e}} & \movesto & \ensuremath{\Varid{e}} \\
\rulename{choose-l} & \ensuremath{\Varid{e}\;\rule[-0.1em]{0.15em}{0.75em}\;\textbf{fail}} & \movesto & \ensuremath{\Varid{e}} \\
\rulename{choose-assoc} & \ensuremath{(\Varid{e}_{\mathrm{1}}\;\rule[-0.1em]{0.15em}{0.75em}\;\Varid{e}_{\mathrm{2}})\;\rule[-0.1em]{0.15em}{0.75em}\;\Varid{e}_{\mathrm{3}}} & \movesto & \ensuremath{\Varid{e}_{\mathrm{1}}\;\rule[-0.1em]{0.15em}{0.75em}\;(\Varid{e}_{\mathrm{2}}\;\rule[-0.1em]{0.15em}{0.75em}\;\Varid{e}_{\mathrm{3}})} \\
\rulename{choose} & \ensuremath{\Conid{SX} [\mskip1.5mu CX [\mskip1.5mu \Varid{e}_{\mathrm{1}}\;\rule[-0.1em]{0.15em}{0.75em}\;\Varid{e}_{\mathrm{2}}\mskip1.5mu]\mskip1.5mu]} & \movesto & \ensuremath{\Conid{SX} [\mskip1.5mu CX [\mskip1.5mu \Varid{e}_{\mathrm{1}}\mskip1.5mu]\;\rule[-0.1em]{0.15em}{0.75em}\;CX [\mskip1.5mu \Varid{e}_{\mathrm{2}}\mskip1.5mu]\mskip1.5mu]} &  \\


\end{array} \\
\\[-0.3em]
\text{\textit{Note}: In the rules marked with a superscript $\alpha$, use $\alpha$-conversion to satisfy the side condition.}
\end{array}$
\caption{\textbf{The \vcname: Rewrite rules to be used for the proof of skew confluence (July 14, 2023)}} \label{fig:rulesOLD}
\Description{A table that exhibits the rewrite rules for the Verse Calculus, groupesd into five categories: Application, Unification, Normalization, Garbage Collection, and Choice}
\end{figure}
}


\newcommand\MC[2]{\multicolumn{#1}{@{}l@{}}{#2}}

\newcommand{\figureSkewProofRewriteRules}{
  \begin{figure}
    $\begin{array}{@{}l}
\begin{array}{@{\hskip0.5em} l @{\hspace{-.5em}} r @{\;}c@{\;} l @{\hskip4em} l r c l@{}}
\rulegroupname{Application:}
\rulename{app-add}          & \ensuremath{\textbf{add}\langle\Varid{k}_{\mathrm{1}},\Varid{k}_{\mathrm{2}}\rangle} & \movesto & \MC{2}{\ensuremath{\Varid{k}_{\mathrm{3}}}} & \MC{3}{\text{where $k_3 = k_1 + k_2$}} \\
\rulename{app-gt}           & \ensuremath{\textbf{gt}\langle\Varid{k}_{\mathrm{1}},\Varid{k}_{\mathrm{2}}\rangle}   & \movesto & \MC{2}{\ensuremath{\Varid{k}_{\mathrm{1}}}} &  \MC{3}{\text{if $k_1 > k_2$}} \\
\rulename{app-gt-fail}      & \ensuremath{\textbf{gt}\langle\Varid{k}_{\mathrm{1}},\Varid{k}_{\mathrm{2}}\rangle}   & \movesto & \MC{2}{\ensuremath{\textbf{fail}}} &  \MC{3}{\text{if $k_1 \leq k_2$}} \\
\rulename{app-beta}\rlap{${}^{\alpha}$} & \ensuremath{(\lambda\Varid{x}.\,\Varid{e})(\Varid{v})} & \movesto & \MC{2}{\ensuremath{\exists\Varid{x}.\,\Varid{x}\hspace{-0.14em}=\hspace{-0.14em}\Varid{v};\,\Varid{e}}}
                                                  & \MC{3}{\text{if $x \not\in \freevars{\ensuremath{\Varid{v}}}$}} \\
\rulename{app-tup}          & \ensuremath{\langle{}\Varid{v}_{\mathrm{0}},\mydots,\Varid{v}_{\Varid{n}}\rangle(\Varid{v})} & \movesto &
                              \MC{5}{\ensuremath{\exists\Varid{x}.\,\Varid{x}\hspace{-0.14em}=\hspace{-0.14em}\Varid{v};(\Varid{x}\hspace{-0.14em}=\hspace{-0.14em}\mathrm{0};\Varid{v}_{\mathrm{0}})\;\rule[-0.1em]{0.15em}{0.75em}\;\mydots\;\rule[-0.1em]{0.15em}{0.75em}\;(\Varid{x}\hspace{-0.14em}=\hspace{-0.14em}\Varid{n};\Varid{v}_{\Varid{n}})}
                              \hspace{0.7em}\text{fresh $\ensuremath{\Varid{x}} \not\in \freevars{\ensuremath{\Varid{v},\Varid{v}_{\mathrm{0}},\mydots,\Varid{v}_{\Varid{n}}}}$}} \hskip-0.7em \\
\rulename{app-tup-0}        & \ensuremath{\langle\rangle(\Varid{v})} & \movesto & \MC{2}{\ensuremath{\textbf{fail}}} \\

    \\[-1.0em]

\rulegroupname{Unification:}
  \rulename{u-lit}   & \ensuremath{\Varid{k}_{\mathrm{1}}\hspace{-0.14em}=\hspace{-0.14em}\Varid{k}_{\mathrm{2}};\,\Varid{e}} & \movesto & \MC{2}{\ensuremath{\Varid{e}}} & \MC{3}{\text{if $k_1=k_2$}} \\
  \rulename{u-tup} & \hspace{-2.0em} \ensuremath{\langle{}\Varid{v}_{\mathrm{1}},\mydots,\Varid{v}_{\Varid{n}}\rangle\hspace{-0.14em}=\hspace{-0.14em}\langle{}\Varid{v}_{1}',\mydots,\Varid{v}_{\Varid{n}}'\rangle;\,\Varid{e}} & \movesto & \MC{2}{\ensuremath{\Varid{v}_{\mathrm{1}}\hspace{-0.14em}=\hspace{-0.14em}\Varid{v}_{1}';\,\mydots;\,\Varid{v}_{\Varid{n}}\hspace{-0.14em}=\hspace{-0.14em}\Varid{v}_{\Varid{n}}';\,\Varid{e}}} \\
  \rulename{u-fail} & \ensuremath{\mathit{hnf}_{\!1}\hspace{-0.14em}=\hspace{-0.14em}\mathit{hnf}_{\!2};\,\Varid{e}} & \movesto & \MC{2}{\ensuremath{\textbf{fail}}} &
       \MC{3}{\begin{array}[t]{@{}l@{}} \text{if \rulename{u-lit}, \rulename{u-tup} do not match and} \\
         \text{\hskip0.6em  neither \ensuremath{\mathit{hnf}_{\!1}} nor \ensuremath{\mathit{hnf}_{\!2}} is a lambda}
         \end{array}} \\
  \rulename{u-occurs} & \ensuremath{\Varid{x}\hspace{-0.14em}=\hspace{-0.14em}\Conid{V} [\mskip1.5mu \Varid{x}\mskip1.5mu];\,\Varid{e}} & \movesto & \MC{2}{\ensuremath{\textbf{fail}}} & \MC{3}{\text{if $\ensuremath{\Conid{V}} \ne \ensuremath{\hole}$}} \\
  \rulename{wrap} & \ensuremath{\Varid{x}\mathrel{=}\mathit{hnf};\,\Varid{e}} & \movesto & \MC{2}{\ensuremath{\Varid{x}\mathrel{=}\mu\Varid{x}.\,\mathit{hnf};\,\Varid{e}}} &
                      \MC{3}{\begin{array}[t]{@{}l@{}} \text{if \rulename{u-occurs} does not match and} \\
         \text{\hskip0.6em $\ensuremath{\Varid{x}} \in \freevars{\ensuremath{\mathit{hnf}}}$}
         \end{array}} \\
  \rulename{unwrap} & \ensuremath{\mu\Varid{x}.\,\mathit{hnf}} & \movesto & \MC{2}{\ensuremath{\subst{\mathit{hnf}}{\Varid{x}}{\mu\Varid{x}.\,\mathit{hnf}}}} \\
  \rulename{subst} & \ensuremath{\Varid{x}\hspace{-0.14em}=\hspace{-0.14em}\Varid{v};\,\Varid{e}} & \movesto & \MC{2}{\ensuremath{\Varid{x}\hspace{-0.14em}=\hspace{-0.14em}\Varid{v};\,\subst{\Varid{e}}{\Varid{x}}{\Varid{v}}}} \hspace*{-0.3em} & \MC{3}{\text{if $\ensuremath{\Varid{x}} \notin \freevars{\ensuremath{\Varid{v}}}$}} \\
  \rulename{hnf-swap} & \ensuremath{\mathit{hnf}\hspace{-0.14em}=\hspace{-0.14em}\Varid{x};\,\Varid{e}}    & \movesto & \MC{2}{\ensuremath{\Varid{x}\hspace{-0.14em}=\hspace{-0.14em}\mathit{hnf};\,\Varid{e}}} \\
  \rulename{seq-swap} & \ensuremath{eq;\,\Varid{x}\hspace{-0.14em}=\hspace{-0.14em}\Varid{v};\,\Varid{e}} & \movesto & \rlap{\ensuremath{\Varid{x}\hspace{-0.14em}=\hspace{-0.14em}\Varid{v};\,eq;\,\Varid{e}}}    & \hskip 2.5em\rlap{\rulename{var-swap}} & \ensuremath{\Varid{y}\hspace{-0.14em}=\hspace{-0.14em}\Varid{x};\,\Varid{e}}    & \movesto & \ensuremath{\Varid{x}\hspace{-0.14em}=\hspace{-0.14em}\Varid{y};\,\Varid{e}}\\

    \\[-1.0em]

\rulegroupname{Elimination:}
\rulename{val-elim} & \ensuremath{\Varid{v};\,\Varid{e}} & \movesto & \MC{2}{\ensuremath{\Varid{e}}} \\
\rulename{exi-elim} & \ensuremath{\exists\Varid{x}.\,\Varid{e}} & \movesto & \MC{2}{\ensuremath{\Varid{e}}} & \MC{3}{\text{if $\ensuremath{\Varid{x}} \notin \freevars{\ensuremath{\Varid{e}}}$}} \\
\rulename{eqn-elim} & \ensuremath{\exists\Varid{x}.\,\Varid{x}\hspace{-0.14em}=\hspace{-0.14em}\Varid{v};\,\Varid{e}} & \movesto & \MC{2}{\ensuremath{\Varid{e}}} & \MC{3}{\text{if $\ensuremath{\Varid{x}} \notin \freevars{\ensuremath{\Varid{v}},\ensuremath{\Varid{e}}}$}} \\
\rulename{fail-elim-eq} & \ensuremath{\Varid{x}\mathrel{=}\textbf{fail};\,\Varid{e}} & \movesto & \MC{2}{\ensuremath{\textbf{fail}}} & \\
\rulename{fail-elim-l} & \ensuremath{\textbf{fail};\,\Varid{e}} & \movesto & \MC{2}{\ensuremath{\textbf{fail}}} & \\
\rulename{fail-elim-r} & \ensuremath{eq;\,\textbf{fail}} & \movesto &\MC{2}{\ensuremath{\textbf{fail}}} & \\

    \\[-1.0em]

\rulegroupname{Normalization:}
\rulename{exi-float-eq}\rlap{${}^{\alpha}$} & \ensuremath{\Varid{v}\mathrel{=}(\exists\Varid{x}.\,\Varid{e});\,\Varid{e'}} & \movesto & \MC{2}{\ensuremath{\exists\Varid{x}.\,\Varid{v}\mathrel{=}\Varid{e};\,\Varid{e'}}} & \MC{3}{\text{if $x \notin \freevars{\ensuremath{\Varid{v}},\ensuremath{\Varid{e'}}}$}} \\
\rulename{exi-float-l}\rlap{${}^{\alpha}$} & \ensuremath{(\exists\Varid{x}.\,\Varid{e});\,\Varid{e'}} & \movesto & \MC{2}{\ensuremath{\exists\Varid{x}.\,\Varid{e};\,\Varid{e'}}} & \MC{3}{\text{if $x \notin \freevars{\ensuremath{\Varid{e'}}}$}} \\
\rulename{exi-float-r}\rlap{${}^{\alpha}$} & \ensuremath{eq;\,(\exists\Varid{x}.\,\Varid{e})} & \movesto & \MC{2}{\ensuremath{\exists\Varid{x}.\,eq;\,\Varid{e}}} & \MC{3}{\text{if $x \notin \freevars{\ensuremath{eq}}$}} \\
\rulename{eqn-float} & \ensuremath{\Varid{v}\hspace{-0.14em}=\hspace{-0.14em}(eq;\,\Varid{e}_{\mathrm{1}});\,\Varid{e}_{\mathrm{2}}} & \movesto & \MC{2}{\ensuremath{eq;\,(\Varid{v}\hspace{-0.14em}=\hspace{-0.14em}\Varid{e}_{\mathrm{1}};\,\Varid{e}_{\mathrm{2}})}} \\
\rulename{seq-assoc} & \ensuremath{(eq;\,\Varid{e}_{\mathrm{1}});\,\Varid{e}_{\mathrm{2}}} & \movesto & \MC{2}{\ensuremath{eq;\,(\Varid{e}_{\mathrm{1}};\,\Varid{e}_{\mathrm{2}})}} \\
\rulename{exi-swap} & \ensuremath{\exists\Varid{x}.\,\exists\Varid{y}.\,\Varid{e}} & \movesto & \MC{2}{\ensuremath{\exists\Varid{y}.\,\exists\Varid{x}.\,\Varid{e}}} \\

    \\[-1.0em]

\rulegroupname{Choice:}
\rulename{one-fail}   & \ensuremath{\textbf{one}\{\textbf{fail}\}} & \movesto & \ensuremath{\textbf{fail}}                   & \rulename{all-fail}   & \ensuremath{\textbf{all}\{\textbf{fail}\}} & \movesto & \ensuremath{\langle{}\rangle} \\
\rulename{one-value}  & \ensuremath{\textbf{one}\{\Varid{v}\}} & \movesto & \ensuremath{\Varid{v}}                       & \rulename{all-value}  & \ensuremath{\textbf{all}\{\Varid{v}\}} & \movesto & \ensuremath{\langle{}\Varid{v}\rangle} \\
\rulename{one-choice} & \ensuremath{\textbf{one}\{\Varid{v}\hspace{0.2em}\mathop{\rule[-0.1em]{0.15em}{0.75em}}\hspace{0.2em}\Varid{e}\}} & \movesto & \ensuremath{\Varid{v}}    & \rulename{all-choice} & \ensuremath{\textbf{all}\{\Varid{v}_{\mathrm{1}}\;\rule[-0.1em]{0.15em}{0.75em}\;\mydots\;\rule[-0.1em]{0.15em}{0.75em}\;\Varid{v}_{\Varid{n}}\}} & \movesto & \ensuremath{\langle{}\Varid{v}_{\mathrm{1}},\mydots,\Varid{v}_{\Varid{n}}\rangle} \\
\rulename{choose-r} & \ensuremath{\textbf{fail}\;\rule[-0.1em]{0.15em}{0.75em}\;\Varid{e}} & \movesto & \ensuremath{\Varid{e}}                & \rulename{choose-l} & \ensuremath{\Varid{e}\;\rule[-0.1em]{0.15em}{0.75em}\;\textbf{fail}} & \movesto & \ensuremath{\Varid{e}} \\
\rulename{choose-assoc} & \ensuremath{(\Varid{e}_{\mathrm{1}}\;\rule[-0.1em]{0.15em}{0.75em}\;\Varid{e}_{\mathrm{2}})\;\rule[-0.1em]{0.15em}{0.75em}\;\Varid{e}_{\mathrm{3}}} & \movesto & \MC{2}{\ensuremath{\Varid{e}_{\mathrm{1}}\;\rule[-0.1em]{0.15em}{0.75em}\;(\Varid{e}_{\mathrm{2}}\;\rule[-0.1em]{0.15em}{0.75em}\;\Varid{e}_{\mathrm{3}})}} \\
\rulename{choose} & \ensuremath{\Conid{SX} [\mskip1.5mu CX [\mskip1.5mu \Varid{e}_{\mathrm{1}}\;\rule[-0.1em]{0.15em}{0.75em}\;\Varid{e}_{\mathrm{2}}\mskip1.5mu]\mskip1.5mu]} & \movesto & \MC{2}{\ensuremath{\Conid{SX} [\mskip1.5mu CX [\mskip1.5mu \Varid{e}_{\mathrm{1}}\mskip1.5mu]\;\rule[-0.1em]{0.15em}{0.75em}\;CX [\mskip1.5mu \Varid{e}_{\mathrm{2}}\mskip1.5mu]\mskip1.5mu]}} &  \\

\end{array} \\

    \\[-1.1em]
       
\text{\textit{Note}: In the rules marked with a superscript $\alpha$, use $\alpha$-conversion to satisfy the side condition.}
\end{array}$
\caption{\textbf{The \vcname: Rewrite rules to be used for the proof of skew confluence (July 14, 2023)}} \label{fig:rules}
\Description{A table that exhibits the rewrite rules for the Verse Calculus, groupesd into five categories: Application, Unification, Normalization, Garbage Collection, and Choice}
\end{figure}
}


\newcommand{\figureSkewProofRewriteRulesNEW}{
  \begin{figure}
    $\begin{array}{@{}l}
\begin{array}{@{\hskip0.5em} l @{\hspace{-.5em}} r @{\;}c@{\;} l @{\hskip4em} l l@{}}
%%\rulegroupname{Application:}
\rulename{app-add}          & \ensuremath{\textbf{add}\langle\Varid{k}_{\mathrm{1}},\Varid{k}_{\mathrm{2}}\rangle} & \movesto & \MC{2}{\ensuremath{\Varid{k}_{\mathrm{3}}}} & \text{where $k_3 = k_1 + k_2$} \\
\rulename{app-gt}           & \ensuremath{\textbf{gt}\langle\Varid{k}_{\mathrm{1}},\Varid{k}_{\mathrm{2}}\rangle}   & \movesto & \MC{2}{\ensuremath{\Varid{k}_{\mathrm{1}}}} &  \text{if $k_1 > k_2$} \\
\rulename{app-gt-fail}      & \ensuremath{\textbf{gt}\langle\Varid{k}_{\mathrm{1}},\Varid{k}_{\mathrm{2}}\rangle}   & \movesto & \MC{2}{\ensuremath{\textbf{fail}}} &  \text{if $k_1 \leq k_2$} \\
\rulename{app-beta}\rlap{${}^{\alpha}$} & \ensuremath{(\lambda\Varid{x}.\,\Varid{e})(\Varid{v})} & \movesto & \ensuremath{\exists\Varid{x}.\,\Varid{x}\hspace{-0.14em}=\hspace{-0.14em}\Varid{v};\,\Varid{e}} & \hbox{\hskip6em}
                                                  & \text{if $x \not\in \freevars{\ensuremath{\Varid{v}}}$} \\
\rulename{app-tup}          & \ensuremath{\langle{}\Varid{v}_{\mathrm{0}},\mydots,\Varid{v}_{\Varid{n}}\rangle(\Varid{v})} & \movesto &
                              \MC{3}{\ensuremath{\exists\Varid{x}.\,\Varid{x}\hspace{-0.14em}=\hspace{-0.14em}\Varid{v};(\Varid{x}\hspace{-0.14em}=\hspace{-0.14em}\mathrm{0};\Varid{v}_{\mathrm{0}})\;\rule[-0.1em]{0.15em}{0.75em}\;\mydots\;\rule[-0.1em]{0.15em}{0.75em}\;(\Varid{x}\hspace{-0.14em}=\hspace{-0.14em}\Varid{n};\Varid{v}_{\Varid{n}})}
                              \hspace{0.7em}\text{fresh $\ensuremath{\Varid{x}} \not\in \freevars{\ensuremath{\Varid{v},\Varid{v}_{\mathrm{0}},\mydots,\Varid{v}_{\Varid{n}}}}$}} \hskip-0.7em \\
\rulename{app-tup-0}        & \ensuremath{\langle\rangle(\Varid{v})} & \movesto & \MC{2}{\ensuremath{\textbf{fail}}} \\

    \\[-0.9em]

%%\rulegroupname{Unification:}
  \rulename{u-lit}   & \ensuremath{\Varid{k}_{\mathrm{1}}\hspace{-0.14em}=\hspace{-0.14em}\Varid{k}_{\mathrm{2}};\,\Varid{e}} & \movesto & \MC{2}{\ensuremath{\Varid{e}}} & \text{if $k_1=k_2$} \\
  \rulename{u-tup} & \hspace{-2.0em} \ensuremath{\langle{}\Varid{v}_{\mathrm{1}},\mydots,\Varid{v}_{\Varid{n}}\rangle\hspace{-0.14em}=\hspace{-0.14em}\langle{}\Varid{v}_{1}',\mydots,\Varid{v}_{\Varid{n}}'\rangle;\,\Varid{e}} & \movesto & \MC{3}{\ensuremath{\Varid{v}_{\mathrm{1}}\hspace{-0.14em}=\hspace{-0.14em}\Varid{v}_{1}';\,\mydots;\,\Varid{v}_{\Varid{n}}\hspace{-0.14em}=\hspace{-0.14em}\Varid{v}_{\Varid{n}}';\,\Varid{e}}} \\
  \rulename{u-fail} & \ensuremath{\mathit{hnf}_{\!1}\hspace{-0.14em}=\hspace{-0.14em}\mathit{hnf}_{\!2};\,\Varid{e}} & \movesto & \ensuremath{\textbf{fail}} &
       \MC{2}{\begin{array}[t]{@{}l@{}} \text{if \rulename{u-lit}, \rulename{u-tup} do not match and} \\
         \text{\hskip0.6em  neither \ensuremath{\mathit{hnf}_{\!1}} nor \ensuremath{\mathit{hnf}_{\!2}} is a lambda}
         \end{array}} \\
  \rulename{u-occurs} & \ensuremath{\Varid{x}\hspace{-0.14em}=\hspace{-0.14em}\Conid{V} [\mskip1.5mu \Varid{x}\mskip1.5mu];\,\Varid{e}} & \movesto & \MC{2}{\ensuremath{\textbf{fail}}} & \text{if $\ensuremath{\Conid{V}} \ne \ensuremath{\hole}$} \\
  \rulename{lam-wrap} & \ensuremath{\Varid{x}\mathrel{=}\Conid{V} [\mskip1.5mu \lambda\Varid{y}.\,\Varid{e}\mskip1.5mu];\,\Varid{e'}} & \movesto & \MC{3}{\ensuremath{\Varid{x}\mathrel{=}\Conid{V} [\mskip1.5mu \lambda\Varid{y}.\,\exists\Varid{x}.\,\Varid{x}\mathrel{=}\Conid{V} [\mskip1.5mu \lambda\Varid{y}.\,\Varid{e}\mskip1.5mu];\,\Varid{e}\mskip1.5mu];\,\Varid{e'}}} \\
              &&&&        \MC{2}{\begin{array}[t]{@{}l@{}} \text{if \rulename{u-occurs} does not match and} \\
         \text{\hskip0.6em $\ensuremath{\Varid{x}} \in \freevars{\ensuremath{\lambda\Varid{y}.\,\Varid{e}}}$}
         \end{array}} \\
%%  \rulename{subst} & |x==v;SE^[x]| & \movesto & \MC{2}{|x==v; SE^[v]|} & \text{if $|x| \notin \freevars{|v|}$} \\
  \rulename{subst} & \ensuremath{\Varid{x}\hspace{-0.14em}=\hspace{-0.14em}\Varid{v};\,\Varid{e}} & \movesto & \MC{2}{\ensuremath{\Varid{x}\hspace{-0.14em}=\hspace{-0.14em}\Varid{v};\,\subst{\Varid{e}}{\Varid{x}}{\Varid{v}}}} & \text{if $\ensuremath{\Varid{x}} \notin \freevars{\ensuremath{\Varid{v}}}$} \\
  \rulename{hnf-swap} & \ensuremath{\mathit{hnf}\hspace{-0.14em}=\hspace{-0.14em}\Varid{x};\,\Varid{e}}    & \movesto & \MC{2}{\ensuremath{\Varid{x}\hspace{-0.14em}=\hspace{-0.14em}\mathit{hnf};\,\Varid{e}}} \\
  \rulename{var-swap} & \ensuremath{\Varid{y}\hspace{-0.14em}=\hspace{-0.14em}\Varid{x};\,\Varid{e}}    & \movesto & \ensuremath{\Varid{x}\hspace{-0.14em}=\hspace{-0.14em}\Varid{y};\,\Varid{e}} \\
  \rulename{seq-swap} & \ensuremath{eq;\,\Varid{x}\hspace{-0.14em}=\hspace{-0.14em}\Varid{v};\,\Varid{e}} & \movesto & \ensuremath{\Varid{x}\hspace{-0.14em}=\hspace{-0.14em}\Varid{v};\,eq;\,\Varid{e}} \\

    \\[-0.9em]

%%\rulegroupname{Elimination:}
\rulename{val-elim} & \ensuremath{\Varid{v};\,\Varid{e}} & \movesto & \MC{2}{\ensuremath{\Varid{e}}} \\
\rulename{exi-elim} & \ensuremath{\exists\Varid{x}.\,\Varid{e}} & \movesto & \MC{2}{\ensuremath{\Varid{e}}} & \text{if $\ensuremath{\Varid{x}} \notin \freevars{\ensuremath{\Varid{e}}}$} \\
\rulename{eqn-elim} & \ensuremath{\exists\Varid{x}.\,\Varid{x}\hspace{-0.14em}=\hspace{-0.14em}\Varid{v};\,\Varid{e}} & \movesto & \MC{2}{\ensuremath{\Varid{e}}} & \text{if $\ensuremath{\Varid{x}} \notin \freevars{\ensuremath{\Varid{v}},\ensuremath{\Varid{e}}}$} \\
\rulename{fail-elim-eq} & \ensuremath{\Varid{x}\mathrel{=}\textbf{fail};\,\Varid{e}} & \movesto & \MC{2}{\ensuremath{\textbf{fail}}} & \\
\rulename{fail-elim-l} & \ensuremath{\textbf{fail};\,\Varid{e}} & \movesto & \MC{2}{\ensuremath{\textbf{fail}}} & \\
\rulename{fail-elim-r} & \ensuremath{eq;\,\textbf{fail}} & \movesto &\MC{2}{\ensuremath{\textbf{fail}}} & \\

    \\[-0.9em]

%%\rulegroupname{Normalization:}
\rulename{exi-float-eq}\rlap{${}^{\alpha}$} & \ensuremath{\Varid{v}\mathrel{=}(\exists\Varid{x}.\,\Varid{e});\,\Varid{e'}} & \movesto & \MC{2}{\ensuremath{\exists\Varid{x}.\,\Varid{v}\mathrel{=}\Varid{e};\,\Varid{e'}}} & \text{if $x \notin \freevars{\ensuremath{\Varid{v}},\ensuremath{\Varid{e'}}}$} \\
\rulename{exi-float-l}\rlap{${}^{\alpha}$} & \ensuremath{(\exists\Varid{x}.\,\Varid{e});\,\Varid{e'}} & \movesto & \MC{2}{\ensuremath{\exists\Varid{x}.\,\Varid{e};\,\Varid{e'}}} & \text{if $x \notin \freevars{\ensuremath{\Varid{e'}}}$} \\
\rulename{exi-float-r}\rlap{${}^{\alpha}$} & \ensuremath{eq;\,(\exists\Varid{x}.\,\Varid{e})} & \movesto & \MC{2}{\ensuremath{\exists\Varid{x}.\,eq;\,\Varid{e}}} & \text{if $x \notin \freevars{\ensuremath{eq}}$} \\
\rulename{eqn-float} & \ensuremath{\Varid{v}\hspace{-0.14em}=\hspace{-0.14em}(eq;\,\Varid{e}_{\mathrm{1}});\,\Varid{e}_{\mathrm{2}}} & \movesto & \MC{2}{\ensuremath{eq;\,(\Varid{v}\hspace{-0.14em}=\hspace{-0.14em}\Varid{e}_{\mathrm{1}};\,\Varid{e}_{\mathrm{2}})}} \\
\rulename{seq-assoc} & \ensuremath{(eq;\,\Varid{e}_{\mathrm{1}});\,\Varid{e}_{\mathrm{2}}} & \movesto & \MC{2}{\ensuremath{eq;\,(\Varid{e}_{\mathrm{1}};\,\Varid{e}_{\mathrm{2}})}} \\
\rulename{exi-swap} & \ensuremath{\exists\Varid{x}.\,\exists\Varid{y}.\,\Varid{e}} & \movesto & \MC{2}{\ensuremath{\exists\Varid{y}.\,\exists\Varid{x}.\,\Varid{e}}} \\

    \\[-0.9em]

%%\rulegroupname{Choice:}
\rulename{one-fail}   & \ensuremath{\textbf{one}\{\textbf{fail}\}} & \movesto & \ensuremath{\textbf{fail}} \\
\rulename{one-value}  & \ensuremath{\textbf{one}\{\Varid{v}\}} & \movesto & \ensuremath{\Varid{v}} \\
\rulename{one-choice} & \ensuremath{\textbf{one}\{\Varid{v}\hspace{0.2em}\mathop{\rule[-0.1em]{0.15em}{0.75em}}\hspace{0.2em}\Varid{e}\}} & \movesto & \ensuremath{\Varid{v}} \\
\rulename{all-fail}   & \ensuremath{\textbf{all}\{\textbf{fail}\}} & \movesto & \ensuremath{\langle{}\rangle} \\
\rulename{all-value}  & \ensuremath{\textbf{all}\{\Varid{v}\}} & \movesto & \ensuremath{\langle{}\Varid{v}\rangle} \\
\rulename{all-choice} & \ensuremath{\textbf{all}\{\Varid{v}_{\mathrm{1}}\;\rule[-0.1em]{0.15em}{0.75em}\;\mydots\;\rule[-0.1em]{0.15em}{0.75em}\;\Varid{v}_{\Varid{n}}\}} & \movesto & \ensuremath{\langle{}\Varid{v}_{\mathrm{1}},\mydots,\Varid{v}_{\Varid{n}}\rangle} \\
\rulename{choose-r} & \ensuremath{\textbf{fail}\;\rule[-0.1em]{0.15em}{0.75em}\;\Varid{e}} & \movesto & \ensuremath{\Varid{e}} \\
\rulename{choose-l} & \ensuremath{\Varid{e}\;\rule[-0.1em]{0.15em}{0.75em}\;\textbf{fail}} & \movesto & \ensuremath{\Varid{e}} \\
\rulename{choose-assoc} & \ensuremath{(\Varid{e}_{\mathrm{1}}\;\rule[-0.1em]{0.15em}{0.75em}\;\Varid{e}_{\mathrm{2}})\;\rule[-0.1em]{0.15em}{0.75em}\;\Varid{e}_{\mathrm{3}}} & \movesto & \MC{2}{\ensuremath{\Varid{e}_{\mathrm{1}}\;\rule[-0.1em]{0.15em}{0.75em}\;(\Varid{e}_{\mathrm{2}}\;\rule[-0.1em]{0.15em}{0.75em}\;\Varid{e}_{\mathrm{3}})}} \\
\rulename{choose} & \ensuremath{\Conid{SX} [\mskip1.5mu CX [\mskip1.5mu \Varid{e}_{\mathrm{1}}\;\rule[-0.1em]{0.15em}{0.75em}\;\Varid{e}_{\mathrm{2}}\mskip1.5mu]\mskip1.5mu]} & \movesto & \MC{2}{\ensuremath{\Conid{SX} [\mskip1.5mu CX [\mskip1.5mu \Varid{e}_{\mathrm{1}}\mskip1.5mu]\;\rule[-0.1em]{0.15em}{0.75em}\;CX [\mskip1.5mu \Varid{e}_{\mathrm{2}}\mskip1.5mu]\mskip1.5mu]}}  \\

\end{array} \\

    \\[-1.1em]
       
\text{\textit{Note}: In the rules marked with a superscript $\alpha$, use $\alpha$-conversion to satisfy the side condition.}
\end{array}$
\caption{\textbf{The \vcname: Rewrite rules to be used for the proof of skew confluence (July 18, 2023)}} \label{fig:rulesNEW}
\Description{A table that exhibits the rewrite rules for the Verse Calculus, groupesd into five categories: Application, Unification, Normalization, Garbage Collection, and Choice}
\end{figure}
}


\newcommand{\figureSkewProofContexts}{
\begin{figure}
    $$
\begin{array}{l}
      \begin{array}{lr@{\;}r@{\;}l}
      % \text{Substitution contexts}  & SE     & ::= & SV \vb |SQ ; e| \vb |eq ; SE| \vb |def x SE|
      %                                                \vb |SE `choice` e| \vb |e `choice` SE| \\
      %                                        && \vb & |apply1N SV v| \vb |apply1N v SV| \vb |one SE| \vb |all SE| \\
      %                        & SQ     & ::= & SE \vb |SV == e|  \vb |v == SE| \\
      %        & SV & ::= & |x|  \vb |k| \vb |op| \vb |arr (SV_1, xdots, SV_n)| \vb |lam x SE| \\
      \text{Value contexts} & V & ::= & \hole \vb \ensuremath{\langle{}\Varid{v}_{\mathrm{1}},\mydots,\Conid{V},\mydots,\Varid{v}_{\Varid{n}}\rangle}\\
      \text{Scope contexts} & SX & ::= & \ensuremath{\textbf{one}\{\Conid{SC}\}} \vb \ensuremath{\textbf{all}\{\Conid{SC}\}} \\
                            & SC & ::= & \ensuremath{\hole} \vb \ensuremath{\Conid{SC}\;\rule[-0.1em]{0.15em}{0.75em}\;\Varid{e}} \vb \ensuremath{\Varid{e}\;\rule[-0.1em]{0.15em}{0.75em}\;\Conid{SC}} \\
      \text{Choice contexts} & CX & ::= & \hole \vb \ensuremath{\Varid{v}\hspace{-0.14em}=\hspace{-0.14em}CX;\,\Varid{e}} \vb \ensuremath{CX;\,\Varid{e}} \vb \ensuremath{\Varid{ceq};\,CX} \vb \ensuremath{\exists\Varid{x}.\,CX} \\
      \text{Choice-free exprs} & ce  & ::=  & \ensuremath{\Varid{v}} \vb \ensuremath{\Varid{ceq};\,\Varid{ce}}
                                     \vb \ensuremath{\textbf{one}\{\Varid{e}\}} \vb \ensuremath{\textbf{all}\{\Varid{e}\}} \vb \ensuremath{\exists\Varid{x}.\,\Varid{ce}}
                                     \vb \ensuremath{\Varid{op}(\Varid{v})} \\
                               & ceq & ::= & \ensuremath{\Varid{ce}} \vb \ensuremath{\Varid{v}\hspace{-0.14em}=\hspace{-0.14em}\Varid{ce}} \\
    \end{array}\\
\end{array}
\vspace{-0.5em}
     $$
    \caption{The syntax of contexts used for the proof of skew confluence}
    \Description{A table that exhibits a BNF grammar that defines sequence contexts, execution contexts, value contexts, scope contexts, and choice-free expressions}
      \label{fig:skew-contexts} \label{fig:skew-choice-free-syntax}
\end{figure}
}


\figureSkewProofRewriteRules 
\figureSkewProofRewriteRulesNEW
\figureSkewProofContexts


 
